\section{Arithmetic parameters and integral splittings}
\label{sec:arithmetic}

Our constructions will use matrices, split flags, and gluing maps over one
ring of integers.  We first record the parameter sets that are dense over
that ring and the finite extension needed to remove determinants.

Fix the rank bound $r\geq 1$ and a number field $K\subset\C$ containing
the prescribed boundary eigenvalues.  Choose distinct primes $p,q>r$,
fix a primitive $pq$-th root of unity $\zeta$, and put
\[
 F=K(\sqrt{2},\zeta),\qquad R=\cO_F.
\]
For a projective curve we take $K=\mathbb Q$.  Throughout the proof,
``integral'' means over the full ring $R$, or over the full ring of
integers of a specified finite extension.  In particular, an invertible
integral matrix belongs to $\GL_d(R)$: its determinant is a unit.

\begin{lemma}\label{lem:integral-parameters}
For every $a,b\geq 0$, the set
$R^a\times(R^\times)^b$ is Zariski dense in
$\C^a\times(\C^\times)^b$.  For every $d\geq 1$, the group
$\GL_d(R)$ is Zariski dense in $\GL_d(\C)$.
\end{lemma}

\begin{proof}
The sets $R$ and $R^\times$ are infinite.  For the unit group,
use the powers of the unit $1+\sqrt{2}$, whose inverse is
$\sqrt{2}-1$.  A polynomial vanishing on a product of infinite subsets
of $\C$ is zero, by induction on the number of variables.  Multiplying
a Laurent polynomial by a monomial gives the same conclusion for the
stated product with multiplicative factors.

For $d\geq 2$, the Zariski closure of $\GL_d(R)$ contains every elementary
one-parameter unipotent subgroup, since its parameters in $R$ are dense.
These subgroups generate $\SL_d(\C)$.  The closure also contains the full
diagonal torus, by the first assertion.  Together these generate
$\GL_d(\C)$.  The case $d=1$ is unit density.
\end{proof}

We will always supply free graded modules as part of integral flag data.
This requirement permits the following elementary splitting argument over
$R$, without any assumption on its ideal class group.

\begin{lemma}\label{lem:free-splittings}
Let $0\to A\to B\to C\to 0$ be an exact sequence of $R$-modules, with
$A$ and $C$ finite free.  Then the sequence splits over $R$.  Its sections
form a nonempty affine space under $\Hom_R(C,A)$, and their integral
points are Zariski dense in the complex space of sections.
Consequently, a finite filtration with specified free graded modules
admits a splitting over $R$ that respects those graded identifications.
\end{lemma}

\begin{proof}
Lift an $R$-basis of $C$ to $B$ and extend linearly to obtain a section.
The difference of any two sections is an arbitrary map from $C$ to $A$.
After choosing bases this is a free matrix module, so density follows
from Lemma~\ref{lem:integral-parameters}.  Apply the splitting argument
successively to the quotients of a filtration.
\end{proof}

In particular, every matrix changing between two such integral splittings
is an $R$-isomorphism; its inverse is integral.  This argument concerns
flags whose graded modules are explicitly free.  We will establish that
property for each flag used below.

\begin{lemma}\label{lem:unit-spectrum}
For every $1\leq n\leq r$, the units
\[
 \alpha_i=\zeta^{i-1},\qquad 1\leq i\leq n,
\]
are distinct, and $\alpha_i-\alpha_j\in R^\times$ whenever $i\ne j$.
\end{lemma}

\begin{proof}
The nonzero integer $i-j$ has absolute value less than both $p$ and $q$,
so it is prime to $pq$.  Thus $\zeta^{i-j}$ is primitive of order $pq$.
The algebraic integer $1-\zeta^{i-j}$ has norm
$\Phi_{pq}(1)=1$ in $\mathbb Q(\zeta)$, and hence is a unit.  Here the
value follows directly by taking $T\to1$ in
\[
 \Phi_{pq}(T)=
 \frac{(T^{pq}-1)(T-1)}{(T^p-1)(T^q-1)}.
\]
Multiplication by the unit $\zeta^{j-1}$ proves the claim.
\end{proof}

The unit differences in Lemma~\ref{lem:unit-spectrum} allow spectral
projectors to be formed over $R$.  The freeness of their images, when
needed, will follow from the free graded lines of the accompanying flag.

\begin{lemma}\label{lem:uniform-field}
The field
\[
 L=F\bigl(\zeta_r,\{u^{1/r}:u\in R^\times\}\bigr)
\]
is a finite extension of $F$.  It contains every $r$-th root of every
unit of $R$, and every such root belongs to $\cO_L^\times$.
\end{lemma}

\begin{proof}
By Dirichlet's unit theorem~\cite[Theorem~5.1]{MilneANT}, choose finitely
many generators
$u_0,\ldots,u_s$ of $R^\times$, including a generator of its finite
torsion subgroup.  Choose $v_i$ with $v_i^r=u_i$.  The finite extension
\[
 F(\zeta_r,v_0,\ldots,v_s)
\]
contains all $r$-th roots of all units: if
$u=\prod_i u_i^{m_i}$ with $m_i\in\Z$, its roots are
$\zeta_r^j\prod_i v_i^{m_i}$.
Both $v_i$ and $v_i^{-1}$ are integral, by the monic equations
$T^r-u_i$ and $T^r-u_i^{-1}$, respectively.  Hence all these roots are
units in $\cO_L$.
\end{proof}

The field $L$ depends only on the fixed initial data.  Later we will take
roots of determinants of matrices over $R$; Lemma~\ref{lem:uniform-field}
handles all those roots simultaneously, regardless of the representation
or irreducible component.
